\documentclass[11pt,a4paper]{article}
\usepackage[margin=25mm]{geometry}
\usepackage{amsmath,amssymb,amsthm,hyperref}
\setlength{\emergencystretch}{2em}
\newtheorem{lemma}{Lemma}
\newtheorem{proposition}[lemma]{Proposition}
\newcommand{\gen}[1]{\langle#1\rangle}
\title{Orbit distance and strict support deficits}
\author{Human Compression project: bounded L044 proof component}
\date{}
\begin{document}\maketitle
Use the standing assumptions of the internal-orbit-block note:
$G\leq A_n$, $n\geq12$, has a transitive connected string
C-marking with every maximal omission intransitive; a unique
eligible interior $i$-edge joins $a\in A$ to $b\in B$;
$L=G_{<i}$ fixes $B$ pointwise and is intransitive on $A$.
The $L$-orbits in $A$ have a common size $k$, their actions
are faithful, and $\beta=a^L$ is exchanged with a distinct
$\beta'$ by $\tau=\rho_{i+1}$. All later generators fix
$\beta$ pointwise, and those of index at least $i+3$ fix
$\beta'$ pointwise as well. These facts have already been
proved completely in that note.

\begin{lemma}[Distance to an active commuting orbit]\label{distance}
If $h>i$ and $\rho_h$ acts nontrivially on $A$, put $d=h-i$.
For $d\geq2$, at least $d$ distinct $L$-orbits occur in $A$.
If the first orbit reached on which $\rho_h$ acts nontrivially
is moved to a different orbit by $\rho_h$, at least $d+1$
distinct orbits occur. In all cases $k\geq i+1$.
\end{lemma}
\begin{proof}
Consider the orbit graph of $L$ using only generators with
indices strictly between $i$ and $h$, starting at $\beta$.
Some reachable orbit has nontrivial $\rho_h$-action. Indeed,
if $\rho_h$ fixed all those orbits pointwise, every higher
indexed generator would fix them pointwise too: it fixes
$\beta$ and commutes with all the generators used in the
paths. Their union would be invariant under $G_i$.
Transitivity of $G_i$ on $A$ would then make it all of $A$,
contrary to the nontrivial $\rho_h$-action.

Take a shortest orbit path to such an orbit, of length $t$,
and let $w$ be its word in distinguished generators. Its
orbits are distinct, and $a^w$ is moved by $\rho_h$:
a nontrivial permutation commuting with the transitive
$L$-action on an invariant orbit has no fixed point; if it
moves that orbit the assertion is immediate.

Every index $j$ strictly between $i$ and $h$ occurs in $w$.
Otherwise split the word into factors with indices less
than $j$ and greater than $j$. They commute. The latter
factor fixes $a$, since all its indices are at least $i+2$.
The former commutes with $\rho_h$, which also fixes $a$.
Thus $a^w$ would be fixed by $\rho_h$, a contradiction.
Consequently $t\geq d-1$, proving the first count.
All earlier orbits on the path are fixed pointwise by
$\rho_h$, by the choice of the first active orbit.
If the final orbit is moved, its image is not one of those
earlier fixed orbits, giving the extra orbit claimed.
Finally the faithful degree-$k$ action of $L$ retains its
independent distinguished family of size $i$; the paid
Whiston bound gives $i\leq k-1$.
\end{proof}

\begin{proposition}[The branch with tail supported on the second orbit]
Let $h>i$ be the largest index acting nontrivially on both
$A$ and $B$, and suppose $h<r-1$ and $G_{>h}$ fixes $A$
pointwise. Put $X=\operatorname{supp}(G_{>h})$ and
$F=n-|X|$. Then
\[
                     X\subseteq B\setminus\{b\},\qquad F\geq2h+1.
\]
If $h\geq i+2$, the inequality is strict. If $F=2h+1$,
then $h=i+1$, $k=i+1$, $A=\beta\cup\beta'$, and the fixed
set is exactly $A\cup\{b\}$. The faithful action of $L$
on either orbit is the full $S_k$.
\end{proposition}
\begin{proof}
The accepted fixed-prefix lemma gives, in $B$, a simple path
from $b$ to a point moved by $\rho_h$, using only earlier
indices and excluding $i$. Its $d=h-i$ or more vertices
are fixed by $G_{>h}$. Together with all of $A$, this gives
\[
                              F\geq |A|+d.               \tag{1}
\]
The support inclusion also follows from fixation of $A$
and of the crossing endpoint $b$.

For $d=1$, there are at least two $L$-orbits, so
$F\geq2k+1\geq2i+3=2h+1$. Equality forces
$k=i+1$, exactly two orbits in $A$, and no fixed vertex
outside $A\cup\{b\}$. Whiston's equality theorem makes
the faithful $L$-action on a single orbit $S_k$.

For $d=2$, the group generated by indices between $i$ and
$h$ is just $\gen\tau$. The first active orbit is $\beta'$.
If $\rho_h$ moves it, Lemma~\ref{distance} gives
\[
                       F\geq3k+2\geq3i+5>2i+5=2h+1.
\]
If it preserves it, the complete internal-block note
applies and gives the stronger $F\geq2h+3$.

For $d\geq3$, Lemma~\ref{distance} and (1) give
\[
 F\geq d(k+1)\geq d(i+2)
      =2h+1+\bigl(i(d-2)-1\bigr)\geq2h+1.
\]
Equality could occur only when $i=1$, $d=3$, $k=2$,
and $A$ consists of exactly three $L$-orbits.
But then $L=\gen{\rho_0}$, and on each of those two-point
orbits $\rho_0$ is a transposition. It fixes $B$ pointwise,
so it would be a product of exactly three transpositions,
contrary to evenness. This excludes equality and completes
the proof in every case.
\end{proof}

\paragraph{Exact scope and attribution.}
This is a new orbit-count proof of the numerical support
part of L044 Proposition~5.23 when the final tail fixes the
first orbit, retaining the source's commuting-block idea
and all hypotheses. The source is Cameron--Fernandes--Leemans--Mixer,
\emph{Highest rank of a polytope for $A_n$}, arXiv:1605.09173v1,
physical pages~32--33. Whole orbit counts replace the need
to infer disjointness of transported point paths. The paid
fixed-prefix, internal-block and Whiston inputs remain
explicit. No finite computation, new formal declaration,
whole-original or question credit is asserted.
The branch where the final tail acts on $A$, the full
coloured equality graph, and subsequent full-IP rigidity
are not proved by this note.
\end{document}
